\section{Introduction}
\label{sec:introduction}

Automated resolution of software issues---from bug reports to feature requests---has emerged as a central application of large language models (LLMs) in software engineering \citep{swebench2024, yang2024sweagent, xia2024agentless, zhang2024autocoderover}. A critical first step in any issue resolution pipeline is \emph{file localization}: identifying which files in a repository are relevant to the issue at hand. The quality of this retrieval step directly constrains downstream patch generation, as irrelevant or missing files lead to incorrect or incomplete fixes.

The dominant approach to file localization is dense retrieval, typically implemented as retrieval-augmented generation (RAG) \citep{lewis2020rag}: code is chunked, embedded into a vector index, and the top-$K$ chunks most similar to the issue description are retrieved. While effective, this approach has two structural limitations. First, it is \emph{expensive to set up}: embedding an entire codebase requires substantial API calls, and the resulting index must be rebuilt whenever the repository state changes. Second, it is \emph{structurally unaware}: embedding similarity captures lexical and semantic relatedness but does not encode the dependency structure---imports, function calls, class inheritance---that developers routinely use to navigate code.

We propose \emph{GraphManager}, a retrieval framework that operates over a graph world model derived from the repository's abstract syntax tree (AST). The graph encodes files, classes, and functions as nodes, with typed edges for structural relations (imports, calls, inheritance, containment). A FAISS vector index over node metadata provides semantic entry points into the graph, while graph traversal supplies structural context that flat retrieval misses.

The framework supports a spectrum of retrieval policies: from a fully deterministic graph scorer that uses no LLM calls at query time, to LLM-guided interactive agents that navigate the graph through function-calling tools. This range enables direct comparison of retrieval quality against cost, and isolates the contribution of graph structure from that of LLM reasoning.

A key design principle is \emph{cost awareness}. We decompose total token cost into three components---setup embedding, runtime LLM, and query embedding---and report all three for every method. This accounting reveals a structural advantage of graph-based retrieval: because the graph index is built from compact structural metadata rather than full code chunks, its setup cost is substantially lower than that of a dense code embedding index. When the same graph serves multiple retrieval tasks against a single repository snapshot, this setup cost is amortized further.

To evaluate this tradeoff rigorously, we introduce a dual-track evaluation protocol:
\begin{enumerate}[leftmargin=1.5em]
  \item \textbf{Strict commit fidelity}: each issue is evaluated at its own \texttt{base\_commit}, requiring per-commit index rebuilds and providing a conservative cost estimate.
  \item \textbf{Same-snapshot amortized}: all issues for a repository are evaluated against a single fixed snapshot, measuring the cost benefit of world-model reuse.
\end{enumerate}

We evaluate seven retrieval methods---including deterministic graph scoring, LLM-guided graph navigation, LLM-guided RAG, and non-LLM RAG baselines---on issues drawn from SWE-bench \citep{swebench2024} across multiple Python repositories.

\paragraph{Contributions.}
\begin{enumerate}[leftmargin=1.5em]
  \item A graph-guided retrieval framework with seven method variants, including a deterministic multi-component scorer that requires no LLM calls at query time.
  \item An amortization-aware dual-track evaluation protocol that separately reports strict commit-fidelity and same-snapshot reuse settings, preventing conflation of retrieval quality with cost structure.
  \item A reproducible cost-quality analysis on SWE-bench showing that graph-based methods achieve competitive file-level F1 at lower total token cost than RAG baselines.
  \item A patching pilot on 100 SWE-bench Verified instances across 9 repositories providing directional evidence that the retrieval cost advantage extends end-to-end to issue resolution.
\end{enumerate}
